%  LaTeX support: latex@mdpi.com 
%  For support, please attach all files needed for compiling as well as the log file, and specify your operating system, LaTeX version, and LaTeX editor.

%=================================================================
\documentclass[environments,article,submit,pdftex,moreauthors]{Definitions/mdpi} 

%=================================================================
%----------
% journal
%----------
% Choose between the following MDPI journals:

% MDPI internal commands
\firstpage{1} 
\makeatletter 
\setcounter{page}{\@firstpage} 
\makeatother
\pubvolume{1}
\issuenum{1}
\articlenumber{0}
\pubyear{2023}
\copyrightyear{2023}
%\externaleditor{Academic Editor: Firstname Lastname}
\datereceived{25 July 2023} 
\daterevised{} 
\dateaccepted{} 
\datepublished{} 
\hreflink{https://doi.org/} % If needed use \linebreak
%\doinum{}

%=================================================================
% Add packages and commands here. The following packages are loaded in our class file: fontenc, inputenc, calc, indentfirst, fancyhdr, graphicx, epstopdf, lastpage, ifthen, lineno, float, amsmath, setspace, enumitem, mathpazo, booktabs, titlesec, etoolbox, tabto, xcolor, soul, multirow, microtype, tikz, totcount, changepage, attrib, upgreek, cleveref, amsthm, hyphenat, natbib, hyperref, footmisc, url, geometry, newfloat, caption

\graphicspath{{figures/}} % path to folder with figures
\usepackage{subcaption}
\usepackage{listings}
\usepackage{xcolor}
\usepackage{gensymb} % for \textdegree
\usepackage{tabularx}

%=================================================================
% Full title of the paper (Capitalized)
\Title{Monitoring Seasonal Fluctuations in Sabkhas of Tunisia by Multispectral Image Analysis Using GRASS GIS}

% MDPI internal command: Title for citation in the left column
\TitleCitation{Monitoring Seasonal Fluctuations in Sabkhas of Tunisia by Multispectral Image Analysis Using GRASS GIS}

% Author Orchid ID: enter ID or remove command
\newcommand{\orcidauthorA}{0000-0002-5759-1089} % Polina Add \orcidA{} behind the author's name

% Authors, for the paper (add full first names)
\Author{Polina Lemenkova $^{1}$*\orcidA{}}

%\longauthorlist{yes}

% MDPI internal command: Authors, for metadata in PDF
\AuthorNames{Polina Lemenkova}

% MDPI internal command: Authors, for citation in the left column
\AuthorCitation{Lemenkova, P.}
% If this is a Chicago style journal: Lastname, Firstname, Firstname Lastname, and Firstname Lastname.

% Affiliations / Addresses (Add [1] after \address if there is only one affiliation.)
\address{%
$^{1}$ \quad Laboratory of Image Synthesis and Analysis (LISA), École polytechnique de Bruxelles (Brussels Faculty of Engineering), Université Libre de Bruxelles (ULB). Building L, Campus du Solbosch, ULB -- LISA CP165/57, Avenue Franklin D. Roosevelt 50, Brussels 1050, Belgium.
}

% Contact information of the corresponding author
\corres{Correspondence: polina.lemenkova@ulb.be; Tel.: +32 471 86 04 59 (P.L.)}

% Current address and/or shared authorship
%\firstnote{Current address: Affiliation 3.} 

% Abstract (Do not insert blank lines, i.e. \\) 
\abstract{This study documents the variations of the saline lakes (a.k.a. sebkhet or sabkha) in Tunisia due to the climate effects. Located in arid environment of north Africa, salt lakes experience changes in volume and extent due to extreme temperatures: during summer heat period, lakes dry out almost entirely and become a crust of salt and minerals due to the evaporated moisture. Lakes in Tunisia are important components of the environment playing a crucial role as biodiversity maintenance, habitats for rare species in nearby oases, rare sources of water, flood and fishery, regulators of temperature, hydrology, and carbon fixation. Remote sensing data are valuable data source to monitoring lakes since variations in the extent of lakes are visible from space. In this study, changes in salinisation and water extent of the salt pans of Tunisia were evaluated using a series of Landsat 8-9 OLI/TIRS images processed with GRASS GIS software. The study area included four salt lakes of north Tunisia: 1) Sebkhet de Sidi el Hani (Sousse Governorate), 2) Sebkha de Moknine (Mahdia Governorate), 3) Sebkhet El Rharra and 4) Sebkhet en Noual (Sfax). The scenes were analysed for a periods 2013--2023. Spatio-temporal changes and their climate-environmental driving forces were analysed. The results were interpreted and the highest changes detected by accuracy assessment, computing the kappa confusion matrices, evaluating the confidence intervals, and plotting the reject probability maps. Multi-temporal changes in land cover classes are reported for each image. The results demonstrated changes in salt lakes were determined for winter/spring/summer months as detected changes in water/land/salt/sand/vegetation areas. This paper contributes to the environmental monitoring of Tunisian landscapes and analysis of climate effects on landscapes of north Africa.}

% Keywords
\keyword{Remote sensing; image processing; Africa; Sahara; salt lakes; desert environment; arid climate; Tunisia; image analysis; Landsat} 

% The fields PACS, MSC, and JEL may be left empty or commented out if not applicable
\PACS{91.10.Da; 91.10.Jf; 91.10.Sp; 91.10.Xa; 96.25.Vt; 91.10.Fc; 95.40.+s; 95.75.Qr; 95.75.Rs; 42.68.Wt}
\MSC{86A30; 86-08; 86A99; 86A04}
\JEL{Y91; Q20; Q24; Q23; Q3; Q01; R11; O44; O13; Q5; Q51; Q55; N57; C6; C61}

\begin{document}

%----------- section ----------->
\section{Introduction}

\subsection{Background}

Remote sensing data analysis is a fundamental issue in environmental studies and ecosystem monitoring. Processing remote sensing data is especially essential for multi-temporal analysis of land cover changes \cite{BOUSSETTA2022102564} and application of environmental geomorphology, which are subject to seasonal changes and shrinkage due to climate effects \cite{JonesMillington}. Climate effects can cause desertification \cite{Pontanier,info14040249} and threaten balance of biodiversity \cite{YAICHEACHOUR2023501}. Therefore, monitoring environmental properties of Saharan landscapes by remote sensing data plays an essential role in analysis and decision making on environmental sustainability, which results in many approaches to developing effective methods of image processing for environmental tasks \cite{Monget,Lemenkova202227,CAMPOS2018211}. 

The key factors in effective approaches to satellite image processing include the type of data and the functionality of software \cite{rs14236017}. Another important factor consists in the quality of image, coverage extent, cloudiness and haze. Major image types widely used for effective environmental monitoring and include Landsat \cite{BOUSSEMA2023185,SMIDA2022104643,jmse11040871}, AVHRR \cite{Bryant}, Google Earth \cite{LI2022252}, MODIS \cite{TOUHAMI2022103804,RHIF2022101596}, and the combinations thereof \cite{REBILLARD1984133}. Examples of very high resolution data include such images as Quickbird \cite{CHOUARI2021443}, SPOT and IKONOS \cite{VELA2008591}. Certainly, although diverse images can be used in various tasks as data sources for environmental monitoring and mapping, the general approach of image analysis relies on using imagery as Landsat or other existing sensor types to extract information through image processing \cite{LiangWang}. 

The effects of climate and seasonal meteorological fluctuations on soil and vegetation setting of the Saharan landscapes are explained by changes in temperature which triggers processes of evaporation, i.e., the transform of liquid water in lakes and ephemeral river flows into gaseous vapour. Such increase of temperature and reduced precipitation observed during heat of summer months lead to lack of surface water which is the primary factor of desertification and aridification of landscapes in Sahara \cite{OSWALD2023369,DODORICO2013326}. At the same time, lakes in Tunisia are important components of the dryland environment playing a crucial role as biodiversity maintenance as habitats for rare African species and birds in nearby oases. Moreover, lakes in Sahara are rare sources of water, flood and fishery, they serve as regulators of temperature, hydrology, and carbon fixation for environmental sustainability. 

%----------- section ----------->
\section{Materials and Methods}

%----------- subsection ----------->
\subsection{Data}

In this study, changes in salinisation and water extent of the depressions of salt pans of Tunisia were evaluated using a series of Landsat 8-9 OLI/TIRS images processed with Geographic Resources Analysis Support System (GRASS) GIS software. Tested images included spectral bands in visible, Short-wave infrared (SWIR)-1, SWIR2 and Near-infrared (NIR) channels obtained from the OLI sensor. 

\begin{figure}[H]
\hspace{50pt}\makebox[0.90\linewidth][r]{%
	\begin{subfigure}[b]{.40\textwidth}
		\centering
			\includegraphics[width=0.87\textwidth]{Fig_04a.jpg}
		\caption{28 February 2017}
	\end{subfigure}%
	\begin{subfigure}[b]{.40\textwidth}
		\centering
		\includegraphics[width=0.87\textwidth]{Fig_04b.jpg}
		\caption{1 April 2017}
	\end{subfigure}%
	\begin{subfigure}[b]{.40\textwidth}
		\centering
		\includegraphics[width=0.87\textwidth]{Fig_04c.jpg}
		\caption{6 July 2017}
	\end{subfigure}%
}\\
\vfill \vspace{1mm}
\hspace{50pt}\makebox[0.90\linewidth][r]{%
	\begin{subfigure}[b]{.40\textwidth}
		\centering
			\includegraphics[width=0.87\textwidth]{Fig_04d.jpg}
		\caption{21 February 2023}
	\end{subfigure}%
	\begin{subfigure}[b]{.40\textwidth}
		\centering
		\includegraphics[width=0.87\textwidth]{Fig_04e.jpg}
		\caption{10 April 2023}
	\end{subfigure}%
	\begin{subfigure}[b]{.40\textwidth}
		\centering
		\includegraphics[width=0.87\textwidth]{Fig_04f.jpg}
		\caption{15 July 2023}
	\end{subfigure}%
}
\caption{Landsat images showing increase of salt (bright cyan colour) from winter to summer months in the Tunisian sabkhas Sebkhet de Sidi el Hani (Sousse Governorate), Sebkha de Moknine (Mahdia Governorate) and Sebkhet El Rharra (Sfax Governorate): \textbf{(a)}: 28 February 2017; \textbf{(b)}: 1 April 2017; \textbf{(c)}: 6 July 2017; \textbf{(d)}: 21 February 2023; \textbf{(e)}: 10 April 2023; \textbf{(f)}: 15 July 2023.}\label{fig04}
\end{figure}

%----------- subsection ----------->
\subsection{Methods}
The Landsat 8-9 OLI/TIRS satellite images were processed with Geographic Resources Analysis Support System (GRASS) Geographic Information System (GIS) software. First, the images were processed by scripts using clustering using 'i.cluster' module and continued with classification by the 'i.maxlik' module. The scenes were analysed for 10 recent years from 2013 to 2023 using the maximum-likelihood discriminant analysis classifier. Afterwards, the spatio-temporal changes of water level and salinisation in target salt lakes in Tunisia and its climate-environmental driving forces were analysed. 

The detection of the saline lakes is possible using spectral properties of images that enable to distinct highly reflective cyan-coloured areas of lakes discriminated against beige-coloured sands and bare land. Thus, spectral properties of the remote sensing data were evaluated by various GRASS GIS modules and scripting techniques with regards to the water/land/salt discrimination, and demonstrated robust approach to image analysis. 

%----------- subsection ----------->
\subsection{Subsection}

%----------- section ----------->
\section{Results}

Changes in salt lakes were determined at different years during recent decade by integrating Landsat images taken in winter/spring/summer months for comparison of periods when water inflow into Saharan lakes is high and low, respectively.  

A time series analysis of the satellite images characterises gradual changes in lakes over calendar periods and provides a robust view in environmental changes. 

Changes in water level in salt lakes detected using remote sensing data can serve for environmental monitoring of Tunisian landscapes and analysis of environmental consequences of climate effects in north Africa.

%----------- section ----------->
\section{Discussion}

%----------- section ----------->
\section{Conclusions}

\subsection{Summary}

GRASS GIS software has been successfully applied to identify different stages in the landscape change in sabkhas of north Tunisia (with a special focus on several lakes -- Sebkhet de Sidi el Hani, Sebkha de Moknine, Sebkhet El Rharra and Sebkhet en Noual). The study also evaluated the effects of climate on seasonal changes in the lacustrine landscapes, as indicated by development of land cover types over the period of 2013 to 2023. Besides, the experiments served as a means of evaluating the climate effects as accelerators of salinisation development in sediments of lakes. The presented results confirm that GRASS GIS can be effectively used to processing and classification of remote sensing imagery using programming approach that ensures automation and precision and accuracy of data processing. The outputs shows the fluctuating behaviour of saline lakes in the north Sahara through the investigated periods as a result of the seasonal effects from temperature acting as triggers to salinisation. Moreover, the correlations between the extent of lakes and effects from heat enabled to detect the influence of climate on environment of north Sahara. 

In this study, we have introduced an advanced scripting cartographic approach of GRASS GIS to evaluating lacustrine properties of soil around the sabkhas in north Tunisia using remote sensing data processing. The main idea was to apply in a single workflow framework both the multi-year and the multi-seasonal images in order to reveal the information on saline lakes behaviour with varying dats on image capture with regard to a sequence of different years and seasons within one year. In such a way, the images were analysed as multi-temporal time series on a double level of variation analysis. 

The results of the image processing approach of the GRASS GIS demonstrated in this paper show that the summer period speeds up soil salinisation and mineralisation in sabkhas of Tunisia, as demonstrated through mapped series of images on several seasons with increased summer heat, and vice versa. Then, using the information received from spectral reflectance of the pixels on the images and heat development in Sahara during summer months we were able to find the highest degree of salinisation for each lake in various years. In such a way, the performed experiments on remote sensing data analysis supported monitoring of fluctuations in water level of salt lakes of Tunisia through the analysed time series analysis of satellite images.

\subsection{Recommendations}

For future similar works on remote sensing data analysis of salt lakes, we provide below the following recommendations and technical suggestions for mapping and image analysis:

\begin{enumerate}
	\item Make the initial study similar to the image processing workflow in this report, but also use other GRASS GIS modules applied for processing the target dates. For example, a useful approach is segmentation by 'i.segment'.
	\item Cloudiness vary a lot in different images. It is necessary to use a cloud-free series of images as an input dataset. While in this study, 10\% of cloud coverage was applied, the modified parameters optimised with different cloud content can be tested and applied with decreased cloud coverage. In this way, the use of 'i.landsat.acca' module of the GRSAS GIS enables Automatic Cloud Cover Assessment (ACCA) for a quick evaluation of image quality and suitability.
	\item Process other dataset in addition to the Landsat, e.g., Advanced Spaceborne Thermal Emission and Reflection Radiometer (ASTER) images. ASTER data can be proprocced using special modules of 'i.aster.toar', similar to the 'i.landsat.toar'. The main advantage of ASTER is a higher resolution--a pixel size of 15 m compared to the 30-m Landsat. 
	\item 
\end{enumerate}

\subsection{Concluding remarks}

This paper reports the results obtained from experimental investigations on satellite images Landsat 8-9 OLI/TIRS collected from the USGS. The tests on image processing were designed to evaluate the changes in several saline lakes of north Tunisia in various seasons in 2013 to 2023 and variations in water extent as indicators of salinisation. Specifically, we evaluated the proposed framework on remote sensing data analysis using GRASS GIS software and a series of the satellite images from the USGS repository. We also proposed future directions in similar works to develop robust scheme on remote sensing data processing tests in similar studies where cartographic approach to image analysis is applied in environmental monitoring projects aimed at evaluation of climate effects on landscape changes. In this work we also discussed the ways in which the scripting approach of cartographic data processing contrasts with traditional GIS methods of mapping that use Graphical User Interface (GUI), and complement to RS data processing for evaluating the connections between the climate effects and salinisation of sabkhas in Tunisia. 

We demonstrated the advantages of the GRASS GIS as a quantitative, robust and advanced cartographic tool for evaluating environmental properties based on image analysis. A workflow has been developed that uses a series of high-resolution Landsat 8-9 OLI/TIRS images with various combinations of GRASS GIS modules, evaluated to analyse the properties of lacustrine landscapes  of north Tunisia as a function of summer heat in Sahara and increase in evaporation during the period of 10 years (2013-2023). The approach of GRASS GIS demonstrated to be a more advanced way of image processing compared to traditional GIS software due to scripts that enables automation of workflow repeatability of the process. Thus, as showed in this study, the use of several module of GRASS GIS ensures multifold manipulations with images, such as in 'i.segment', 'i.maxlik' and other sub-tests that where applied for various steps of multispectral image analysis. 

A scripting approach of the GRASS GIS enabled us to map and analyse the rate of salinisation in sabkhas of Tunisia in real time in various seasons from 2013 to 2023 and to obtain the information from the processed images representing the land cover types in the surroundings. In turn, remote sensing datasets enabled to model the behavior of landscapes in spring, summer and winter periods in north Tunisia which was used to compare the climate effects from rise in temperature on the evaporation of water in salt lakes, as directly related to heat increase during summer period in Sahara. The results demonstrated variations in water level in several sabkha and show technical effectiveness of the GRASS GIS in cartographic processing of the remote sensing data for environmental tasks. Correlations between the increase in heat in various years and high salinisation of sabkhas were detected. The experiments on image processing reported in this article have investigated the salinisation behaviour in saline lakes with thermal variations using analysis of Landsat bands in multispectral channels. 

Due to the effectiveness and applicability of the performed experiments on image analysis by GRASS GIS -- flexibility of scripting approach and a wide variety of modules tuned for diverse tasks of image processing using its modular approach -- we expect the proposed method to be able to deal with other satellite imagery that cover diverse regions of north Africa to study development of lacustrine landscapes with effects from climate and fluctuations in temperature on Saharan ecosystems. Various sabkhas of north Tunisia with diverse surrounding soil types and salinisation/mineralisation degree were evaluated to generalise the proposed workflow of salt lake monitoring to a higher dimensional workload in the larger lakes and different world ecosystems with similar salt lakes. 

As demonstrated in this study, the remote sensing data from the USGS EarthExplorer repository can be used to analyse the variations in salt lakes. It enables to derive information using data on spectral reflectance and brightness of pixels which strongly correlates with the distribution of various land cover types. In such a way, remote sensing data processing enables to detect patches in complex diversified lacustrine landscapes of north Sahara characterised with heterogeneous environment in contrasting climate of desert. This paper contributes to the environmental monitoring of Sahara, analysis of salt lake variations and regional studies of Tunisia, north Africa.

%----------- section ----------->

\funding{The publication was funded by the Editorial Office of \emph{Environments}, Multidisciplinary Digital Publishing Institute (MDPI), by providing 100\% discount for the APC of this manuscript.}

\institutionalreview{Not applicable.}

\informedconsent{Not applicable.}

\dataavailability{Not applicable} 

\acknowledgments{The authors thank the reviewers for reading and review of this manuscript.}

\conflictsofinterest{The authors declare no conflict of interest.} 

\abbreviations{Abbreviations}{
The following abbreviations are used in this manuscript:\\

\noindent 
\begin{tabular}{@{}ll}
ACCA & Automatic Cloud Cover Assessment \\
ASTER & Advanced Spaceborne Thermal Emission and Reflection Radiometer \\
GMT & Generic Mapping Tools \\
GIS & Geographic Information System \\
GUI & Graphical User Interface \\
GRASS & Geographic Resources Analysis Support System \\
Landsat OLI/TIRS & Landsat Operational Land Imager and Thermal Infrared Sensor \\
MODIS & Moderate Resolution Imaging Spectroradiometer \\ 
NIR & Near-infrared \\
SPOT & Satellite Pour l’Observation de la Terre \\
SWIR & Short-wave infrared (SWIR)-1 \\
USGS & United States Geological Survey \\
\end{tabular}
}

%%%%%%%%%%%%%%%%%%%%%%%%%%%%%%%%%%%%%%%%%%
%% Optional
\appendixtitles{no} % Leave argument "no" if all appendix headings stay EMPTY (then no dot is printed after "Appendix A"). If the appendix sections contain a heading then change the argument to "yes".
\appendixstart
\appendix
\section[\appendixname~\thesection]{}
\subsection[\appendixname~\thesubsection]{}
The appendix 

\section[\appendixname~\thesection]{}
All appendix sections must be cited in the main text. In the appendices, Figures, Tables, etc. should be labeled, starting with ``A''---e.g., Figure A1, Figure A2, etc.

%%%%%%%%%%%%%%%%%%%%%%%%%%%%%%%%%%%%%%%%%%
\begin{adjustwidth}{-\extralength}{0cm}
%\printendnotes[custom] % Un-comment to print a list of endnotes

\reftitle{References}
%\nocite{*}

%=====================================
% References, variant A: external bibliography
%=====================================
\bibliography{Tunisia.bib}
%\nocite{*}

%%%%%%%%%%%%%%%%%%%%%%%%%%%%%%%%%%%%%%%%%%
\end{adjustwidth}
\end{document}
